\begin{theorem}[Proved: cyclic three-object support through seventeen]
\label{CP2:A-strong-three-seventeen}
A three-object category with cyclic support and at most
seventeen arrows cannot have type $\CP^2$.
\end{theorem}
\begin{proof}
A disconnected category is excluded by connectedness of
the target. An object-retract deletion saves at least
three arrows. If its two remaining objects are strongly
connected, Proposition~\ref{CP2:A-baseline-strong-two}
applies; otherwise their arrow count is at most fourteen
and the complete two-object baseline applies. Assume
there is no object retract.

If all three objects are strongly connected, their
identities are outside the categorical minimal ideal.
Unless all local row counts or all local column counts
are one, both total counts are at least four and the
category has at least $4\cdot4+3=19$ arrows. The exceptional
one-row or one-column case has finite-group homotopy type
by Lemma~\ref{CP2:A-high-Rees}, and is also excluded.

It remains that a proper strong pair $S$ is connected
to one corner monoid $N$. By opposition take $S$ to be
the source component. Let $q=|S|$, $m=|N|$, and let $h$
be the total number of its outgoing arrows. Each of the
two stalks is nonempty, so
\[
 q+m+h\le17,\qquad h\ge2.
\]
If $h=2$, the collage is the full categorical join of
$S$ and $N$ and has zero positive rational products.
We may suppose $h\ge3$. If both $H^2(BS;\mathbb Q)$
and $H^2(BN;\mathbb Q)$ vanish, the two-block lemma
already excludes the target.

Suppose first $H^2(BS;\mathbb Q)\ne0$. Its total minimal
row and column counts $r,c$ are at least three; otherwise
it has finite-group type. Since $q\le13$, the minimal
group is trivial and $rc+2\le13$ forces $r=c=3$.
In particular $q\ge11$, every local row count is at most
two, and $m\le3$. If $m\le2$, the trivial right rational
$N$-module is projective. The objectwise binary projective
cover, Corollary~\ref{CP2:A-baseline-two-row-instances},
kills all second products on this collage.

If $m=3$, the budget forces $q=11$ and $h=3$. The strong
pair then consists exactly of its nine-element categorical
minimal band and its two identities. Choose the object
$b$ whose local row count is one. Every incoming arrow
from its other object lies in the minimal ideal and is
fixed by a local minimal idempotent $e_b$. The only unit
of its corner is its identity, since all its other elements
lie in that ideal. There are no external entrances into
$S$. Choose minimal cross arrows $u,v$ with $vu=e_b$.
All hypotheses of the one-sided object-deletion lemma
hold. Deleting $b$ gives an actual homotopy equivalence
to a two-object category with fewer than seventeen arrows,
contradicting the complete two-object baseline.

Finally suppose $H^2(BS;\mathbb Q)=0$ and
$H^2(BN;\mathbb Q)\ne0$. Then $m\ge5$. If either outgoing
stalk is a singleton, its arrow is fixed by every $N$
actor. Corollary~\ref{CP2:A-fixed-arrow}, with source
block $S$, kills all second products immediately.
Thus both stalks have at least two arrows: $h\ge4$ and
$q\le8$. The minimal group of $S$ is trivial, since a
nontrivial group would require at least
$2\cdot2\cdot2+2=10$ internal arrows. With two identities
outside the ideal, $rc+2\le8$ shows that one total row
or column count is two and the other is at most three.

If $r=2$, there is one minimal row per object. The actual
one-row source-component collapse replaces $S$ by a trivial
object and gives an ordinary left-act cone over $N$.
Here $m\le7$, since $q\ge6$ and $h\ge4$. Its rational
homology is bounded under the proposed target, and the
bounded-nine act theorem excludes it. If $r>2$, necessarily
$r=3$, $c=2$, $q=8$, $m=5$, and $h=4$. The pair consists
of its six-element minimal band and its two identities.
Its one-row object again has all incoming arrows fixed
by its minimal idempotent and no nontrivial units. There
are no external entrances. The same one-sided deletion
gives a two-object category with at most sixteen arrows,
which is excluded. This exhausts every cyclic support.
\end{proof}
